\documentclass[11pt]{article}
\usepackage[margin=1in]{geometry}
\usepackage{amsmath, amssymb}
\usepackage{booktabs}
\usepackage{listings}
\usepackage{xcolor}
\usepackage{hyperref}
\usepackage{graphicx}

% MATLAB code styling
\lstdefinestyle{matlab}{
    language=Matlab,
    basicstyle=\ttfamily\small,
    keywordstyle=\color{blue}\bfseries,
    commentstyle=\color{green!60!black},
    stringstyle=\color{orange},
    numbers=left,
    numberstyle=\tiny\color{gray},
    numbersep=5pt,
    frame=single,
    breaklines=true,
    breakatwhitespace=true,
    tabsize=4,
    showstringspaces=false,
    backgroundcolor=\color{gray!10}
}

\lstset{style=matlab}

\title{MPC Cart-Pole Swing-Up Controller\\[0.5em]\large A Step-by-Step Code Explanation}
\author{}
\date{January 2026}

\begin{document}

\maketitle

\tableofcontents
\newpage

%% ============================================================================
\section{Introduction to Model Predictive Control}
%% ============================================================================

Model Predictive Control (MPC) is a control strategy that solves an optimization problem at each time step to determine the best control action. The key idea is:

\begin{enumerate}
    \item \textbf{Predict}: Simulate the system forward over a finite horizon using a model
    \item \textbf{Optimize}: Find the control sequence that minimizes a cost function
    \item \textbf{Apply}: Apply only the \emph{first} control action
    \item \textbf{Repeat}: Shift the horizon forward and repeat
\end{enumerate}

This ``receding horizon'' approach allows MPC to handle constraints, nonlinear dynamics, and changing conditions naturally.

\subsection{The Cart-Pole Problem}

The cart-pole (inverted pendulum on a cart) is a classic control problem. The goal is to:
\begin{itemize}
    \item Start with the pendulum hanging down ($\theta = 0$)
    \item Swing the pendulum up to the inverted position ($\theta = \pi$)
    \item Stabilize both the pendulum angle and cart position
\end{itemize}

The state vector is:
\[
\mathbf{x} = \begin{bmatrix} q_1 \\ q_2 \\ \dot{q}_1 \\ \dot{q}_2 \end{bmatrix} = \begin{bmatrix} \text{cart position} \\ \text{pendulum angle} \\ \text{cart velocity} \\ \text{angular velocity} \end{bmatrix}
\]

The control input $u$ is the horizontal force applied to the cart.

%% ============================================================================
\section{System Parameters}
%% ============================================================================

The physical parameters of the cart-pole system are defined in a structure:

\begin{lstlisting}
%% System Parameters
params = struct();
params.m1 = 5.0;    % Cart mass (kg)
params.m2 = 1.0;    % Pendulum mass (kg)
params.L = 2.0;     % Pendulum length (m)
params.g = 9.81;    % Gravity (m/s^2)
\end{lstlisting}

These values represent a relatively heavy cart (5 kg) with a lighter pendulum bob (1 kg) on a 2-meter rod. The mass ratio $m_2/m_1 = 0.2$ means the cart is heavy enough to swing the pendulum without excessive motion.

%% ============================================================================
\section{MPC Parameters}
%% ============================================================================

The MPC controller has several tuning parameters:

\begin{lstlisting}
%% MPC Parameters
mpc = struct();
mpc.N = 15;         % Prediction horizon (steps)
mpc.dt = 0.05;      % Time step (20 Hz control)
mpc.u_max = 1000;   % Max force (N)

% Cost weights
mpc.Q = diag([10, 100, 1, 10]);  % [pos, angle, vel_pos, vel_ang]
mpc.R = 0.01;                    % Control weight
mpc.Q_terminal = 2 * mpc.Q;      % Terminal cost weight
\end{lstlisting}

\subsection{Horizon and Timing}

\begin{itemize}
    \item \textbf{N = 15}: The MPC looks 15 steps into the future
    \item \textbf{dt = 0.05}: Each step is 50 ms, so the horizon spans $15 \times 0.05 = 0.75$ seconds
    \item \textbf{Control rate}: $1/0.05 = 20$ Hz
\end{itemize}

A longer horizon allows the MPC to ``see'' further ahead and plan better swing-up trajectories, but increases computation time.

\subsection{Cost Weights}

The cost function penalizes deviations from the target state and control effort:
\[
J = \sum_{k=0}^{N-1} \left[ (\mathbf{x}_k - \mathbf{x}_{\text{ref}})^T \mathbf{Q} (\mathbf{x}_k - \mathbf{x}_{\text{ref}}) + R \, u_k^2 \right] + (\mathbf{x}_N - \mathbf{x}_{\text{ref}})^T \mathbf{Q}_f (\mathbf{x}_N - \mathbf{x}_{\text{ref}})
\]

The weight matrix $\mathbf{Q} = \text{diag}([10, 100, 1, 10])$ means:
\begin{itemize}
    \item Position error weight: 10
    \item \textbf{Angle error weight: 100} (highest priority)
    \item Velocity error weight: 1
    \item Angular velocity error weight: 10
\end{itemize}

The angle weight being 10$\times$ larger than position means the MPC prioritizes getting the pendulum upright over keeping the cart at the target position.

%% ============================================================================
\section{Simulation Setup}
%% ============================================================================

\begin{lstlisting}
%% Simulation Setup
T_sim = 16.0;       % Total simulation time
t_vec = 0:mpc.dt:T_sim;
N_sim = length(t_vec);

% Initial and target states
x0 = [0; 0; 0; 0];              % Start: cart at 0, pendulum down, at rest
x_target = [3; pi; 0; 0];       % Goal: cart at 3m, pendulum up, at rest
\end{lstlisting}

\begin{itemize}
    \item \textbf{Initial state}: Cart at origin, pendulum hanging straight down ($\theta = 0$), everything at rest
    \item \textbf{Target state}: Cart at $x = 3$ m, pendulum inverted ($\theta = \pi$), at rest
\end{itemize}

The challenge is that the pendulum starts in the \emph{stable} equilibrium (down) and must reach the \emph{unstable} equilibrium (up).

%% ============================================================================
\section{The Main MPC Loop}
%% ============================================================================

The core of MPC is the receding horizon loop:

\begin{lstlisting}
%% Main MPC Loop
for k = 1:N_sim-1
    tic;

    %% Solve MPC Problem
    [u_opt, x_pred, cost] = solve_mpc_step(x_current, x_target, params, mpc);

    solve_time(k) = toc;

    %% Apply First Control Action
    u_apply = u_opt(1);
    u_history(k) = u_apply;
    cost_history(k) = cost;

    %% Simulate Real System (with disturbance)
    if k > 100  % Add disturbance after 5 seconds
        disturbance = [0; 0; 0; 0.2*randn()];
    else
        disturbance = zeros(4,1);
    end

    x_next = simulate_one_step(x_current, u_apply, params, mpc.dt, disturbance);

    %% Update State
    x_current = x_next;
    x_history(:,k+1) = x_current;
end
\end{lstlisting}

\subsection{Loop Structure}

At each time step $k$:

\begin{enumerate}
    \item \textbf{Solve}: Call \texttt{solve\_mpc\_step} to find the optimal control sequence $\mathbf{u}^* = [u_0^*, u_1^*, \ldots, u_{N-1}^*]$

    \item \textbf{Apply first}: Only apply $u_0^*$ to the real system (this is the key MPC idea)

    \item \textbf{Simulate}: Advance the real system one time step using RK4 integration

    \item \textbf{Add disturbance}: After 5 seconds, random torque disturbances test the controller's robustness

    \item \textbf{Repeat}: The loop continues with the new state
\end{enumerate}

\subsection{Why Only Apply the First Control?}

The MPC computes a sequence of $N$ controls, but only applies the first one because:
\begin{itemize}
    \item The prediction model isn't perfect
    \item Disturbances affect the real system
    \item Re-solving at each step provides feedback correction
\end{itemize}

%% ============================================================================
\section{The MPC Optimization Function}
%% ============================================================================

\begin{lstlisting}
function [u_opt, x_pred, cost] = solve_mpc_step(x0, x_ref, params, mpc)
%% Solve single MPC optimization problem

% Decision variables: control sequence u = [u(0), u(1), ..., u(N-1)]
n_vars = mpc.N;

% Initial guess (warm start with previous solution if available)
persistent u_prev
if isempty(u_prev)
    u_guess = zeros(n_vars, 1);
else
    % Shift previous solution and add zero at end
    u_guess = [u_prev(2:end); 0];
end

% Bounds on control
lb = -mpc.u_max * ones(n_vars, 1);
ub =  mpc.u_max * ones(n_vars, 1);

% Optimization options (fast for real-time)
options = optimoptions('fmincon', ...
    'Algorithm', 'sqp', ...
    'Display', 'off', ...
    'MaxIterations', 50, ...
    'MaxFunctionEvaluations', 500, ...
    'OptimalityTolerance', 1e-4, ...
    'ConstraintTolerance', 1e-4);

% Solve optimization
[u_opt, cost] = fmincon(@(u) mpc_cost(u, x0, x_ref, params, mpc), ...
                        u_guess, [], [], [], [], lb, ub, ...
                        @(u) mpc_constraints(u, x0, params, mpc), options);

% Store for warm start next time
u_prev = u_opt;
\end{lstlisting}

\subsection{Decision Variables}

The optimizer searches over the control sequence:
\[
\mathbf{u} = [u_0, u_1, \ldots, u_{N-1}]^T \in \mathbb{R}^N
\]

This is called ``control parameterization'' or ``single shooting'' -- the states are computed by simulating forward from $\mathbf{x}_0$ using the control sequence.

\subsection{Warm Starting}

The \texttt{persistent} variable \texttt{u\_prev} stores the previous solution. The next optimization is initialized with:
\[
\mathbf{u}_{\text{guess}} = [u_1^{\text{prev}}, u_2^{\text{prev}}, \ldots, u_{N-1}^{\text{prev}}, 0]
\]

This ``shift and append'' warm start helps because:
\begin{itemize}
    \item Adjacent MPC problems are similar
    \item The shifted solution is usually close to optimal
    \item Reduces solve time significantly
\end{itemize}

\subsection{Optimization Settings}

\begin{itemize}
    \item \textbf{Algorithm: SQP} (Sequential Quadratic Programming) -- good for smooth nonlinear problems
    \item \textbf{MaxIterations: 50} -- limits computation time for real-time feasibility
    \item \textbf{Tolerances: $10^{-4}$} -- loose tolerances trade accuracy for speed
\end{itemize}

%% ============================================================================
\section{The Cost Function}
%% ============================================================================

\begin{lstlisting}
function J = mpc_cost(u, x0, x_ref, params, mpc)
%% MPC Cost Function

% Predict trajectory
x_traj = predict_trajectory(x0, u, params, mpc);

J = 0;

% Running cost
for i = 1:mpc.N-1
    % State cost
    e = x_traj(:,i) - x_ref;
    J = J + e' * mpc.Q * e;

    % Control cost
    J = J + mpc.R * u(i)^2;
end

% Terminal cost
e_terminal = x_traj(:,mpc.N) - x_ref;
J = J + e_terminal' * mpc.Q_terminal * e_terminal;
end
\end{lstlisting}

\subsection{Cost Structure}

The total cost has three parts:

\textbf{1. Running State Cost:}
\[
\sum_{k=0}^{N-2} (\mathbf{x}_k - \mathbf{x}_{\text{ref}})^T \mathbf{Q} (\mathbf{x}_k - \mathbf{x}_{\text{ref}})
\]

\textbf{2. Running Control Cost:}
\[
\sum_{k=0}^{N-2} R \, u_k^2
\]

\textbf{3. Terminal Cost:}
\[
(\mathbf{x}_{N-1} - \mathbf{x}_{\text{ref}})^T \mathbf{Q}_f (\mathbf{x}_{N-1} - \mathbf{x}_{\text{ref}})
\]

The terminal cost weight $\mathbf{Q}_f = 2\mathbf{Q}$ encourages the trajectory to end near the target, partially compensating for the finite horizon.

\subsection{Why Quadratic Costs?}

Quadratic costs are standard because:
\begin{itemize}
    \item They penalize large errors more than small errors
    \item The optimization landscape is smooth (good for gradient-based solvers)
    \item They connect to LQR theory for linear systems
\end{itemize}

%% ============================================================================
\section{Trajectory Prediction}
%% ============================================================================

\begin{lstlisting}
function x_traj = predict_trajectory(x0, u, params, mpc)
%% Predict trajectory using control sequence

x_traj = zeros(4, mpc.N);
x_traj(:,1) = x0;

for i = 1:mpc.N-1
    % Simulate one step forward
    x_traj(:,i+1) = rk4_step(@(t,x) cart_pole_dynamics(t, x, u(i), params), ...
                             0, x_traj(:,i), mpc.dt);
end
end
\end{lstlisting}

This function ``rolls out'' the dynamics: starting from $\mathbf{x}_0$, it simulates forward using the control sequence to produce the predicted state trajectory:
\[
\mathbf{x}_0 \xrightarrow{u_0} \mathbf{x}_1 \xrightarrow{u_1} \mathbf{x}_2 \xrightarrow{u_2} \cdots \xrightarrow{u_{N-1}} \mathbf{x}_N
\]

%% ============================================================================
\section{Cart-Pole Dynamics}
%% ============================================================================

\begin{lstlisting}
function xdot = cart_pole_dynamics(t, x, u, params)
%% Cart-pole dynamics

q1 = x(1);      % Cart position
q2 = x(2);      % Pendulum angle
q1dot = x(3);   % Cart velocity
q2dot = x(4);   % Pendulum angular velocity

% Extract parameters
m1 = params.m1; m2 = params.m2; L = params.L; g = params.g;

% Compute accelerations
q1ddot = cart_accel(q1, q2, q1dot, q2dot, u, m1, m2, L, g);
q2ddot = pendulum_accel(q1, q2, q1dot, q2dot, u, m1, m2, L, g);

% State derivative
xdot = [q1dot; q2dot; q1ddot; q2ddot];
end
\end{lstlisting}

\subsection{Equations of Motion}

The cart-pole dynamics are derived from the Euler-Lagrange equations. The accelerations are:

\textbf{Cart acceleration:}
\begin{lstlisting}
function q1ddot = cart_accel(q1, q2, q1dot, q2dot, u, m1, m2, L, g)
    den = m1 + m2 - m2*cos(q2)^2;
    num = u + m2*sin(q2)*(L*q2dot^2 + g*cos(q2));
    q1ddot = num / den;
end
\end{lstlisting}

\[
\ddot{q}_1 = \frac{u + m_2 \sin(q_2) \left( L \dot{q}_2^2 + g \cos(q_2) \right)}{m_1 + m_2 - m_2 \cos^2(q_2)}
\]

\textbf{Pendulum acceleration:}
\begin{lstlisting}
function q2ddot = pendulum_accel(q1, q2, q1dot, q2dot, u, m1, m2, L, g)
    den = L*(m1 + m2 - m2*cos(q2)^2);
    num = -u*cos(q2) - m2*L*q2dot^2*cos(q2)*sin(q2) - (m1+m2)*g*sin(q2);
    q2ddot = num / den;
end
\end{lstlisting}

\[
\ddot{q}_2 = \frac{-u \cos(q_2) - m_2 L \dot{q}_2^2 \cos(q_2) \sin(q_2) - (m_1 + m_2) g \sin(q_2)}{L (m_1 + m_2 - m_2 \cos^2(q_2))}
\]

These equations capture the coupling between cart and pendulum motion.

%% ============================================================================
\section{Numerical Integration (RK4)}
%% ============================================================================

\begin{lstlisting}
function x_next = rk4_step(f, t, x, dt)
%% Single RK4 integration step

k1 = f(t, x);
k2 = f(t + dt/2, x + dt*k1/2);
k3 = f(t + dt/2, x + dt*k2/2);
k4 = f(t + dt, x + dt*k3);

x_next = x + (dt/6) * (k1 + 2*k2 + 2*k3 + k4);
end
\end{lstlisting}

The 4th-order Runge-Kutta method provides accurate integration:
\[
\mathbf{x}_{k+1} = \mathbf{x}_k + \frac{\Delta t}{6} \left( \mathbf{k}_1 + 2\mathbf{k}_2 + 2\mathbf{k}_3 + \mathbf{k}_4 \right)
\]

where:
\begin{align*}
\mathbf{k}_1 &= f(t_k, \mathbf{x}_k) \\
\mathbf{k}_2 &= f(t_k + \tfrac{\Delta t}{2}, \mathbf{x}_k + \tfrac{\Delta t}{2} \mathbf{k}_1) \\
\mathbf{k}_3 &= f(t_k + \tfrac{\Delta t}{2}, \mathbf{x}_k + \tfrac{\Delta t}{2} \mathbf{k}_2) \\
\mathbf{k}_4 &= f(t_k + \Delta t, \mathbf{x}_k + \Delta t \, \mathbf{k}_3)
\end{align*}

RK4 has local truncation error $O(\Delta t^5)$, making it accurate for the 50 ms time step used here.

%% ============================================================================
\section{Constraints}
%% ============================================================================

\begin{lstlisting}
function [c, ceq] = mpc_constraints(u, x0, params, mpc)
%% MPC Constraints

c = [];   % No inequality constraints
ceq = []; % No equality constraints (dynamics handled in prediction)

% Example of how to add state constraints:
% x_traj = predict_trajectory(x0, u, params, mpc);
% c = [x_traj(1,:)' - 5;    % position < 5
%     -x_traj(1,:)' - 5];   % position > -5
end
\end{lstlisting}

Currently, the only constraints are the control bounds ($|u| \leq 1000$ N), which are handled by \texttt{fmincon}'s \texttt{lb} and \texttt{ub} parameters.

State constraints (e.g., cart position limits) could be added as inequality constraints \texttt{c} $\leq 0$.

%% ============================================================================
\section{Common Issues and Tuning}
%% ============================================================================

\subsection{Swing-Up Challenges}

The swing-up problem is difficult because:
\begin{enumerate}
    \item \textbf{Local minima}: The pendulum hanging down is a stable equilibrium. If position weight is too high, the MPC won't move the cart enough to swing up.

    \item \textbf{Short horizon}: If the horizon is too short, the MPC can't ``see'' that swinging will eventually lead to the upright position.

    \item \textbf{Energy pumping}: Swing-up requires coordinated back-and-forth motion to pump energy into the pendulum.
\end{enumerate}

\subsection{Tuning Guidelines}

\begin{tabular}{lp{8cm}}
\toprule
Parameter & Effect \\
\midrule
\texttt{mpc.N} (horizon) & Longer $\rightarrow$ better planning but slower solve \\
\texttt{mpc.Q(1,1)} (position) & Higher $\rightarrow$ tighter position control but may prevent swing-up \\
\texttt{mpc.Q(2,2)} (angle) & Higher $\rightarrow$ prioritizes getting pendulum upright \\
\texttt{mpc.R} (control) & Higher $\rightarrow$ smoother but slower control \\
\texttt{mpc.Q\_terminal} & Higher $\rightarrow$ encourages reaching target at horizon end \\
\bottomrule
\end{tabular}

\subsection{Two-Phase Approach}

A robust solution is to use different controllers for different phases:
\begin{enumerate}
    \item \textbf{Swing-up phase}: Low position weight, focus on getting pendulum near upright
    \item \textbf{Stabilization phase}: Once $|q_2 - \pi| < \epsilon$, switch to higher position weight
\end{enumerate}

%% ============================================================================
\section{Summary}
%% ============================================================================

The MPC cart-pole controller works by:
\begin{enumerate}
    \item Measuring the current state $\mathbf{x}_k$
    \item Solving an optimization problem over the control sequence $[u_0, \ldots, u_{N-1}]$
    \item Applying only $u_0$ to the real system
    \item Repeating at 20 Hz
\end{enumerate}

Key components:
\begin{itemize}
    \item \textbf{Cost function}: Quadratic penalties on state error and control effort
    \item \textbf{Dynamics}: Euler-Lagrange equations for cart-pole
    \item \textbf{Integration}: RK4 for trajectory prediction
    \item \textbf{Optimization}: SQP via \texttt{fmincon} with warm starting
\end{itemize}

The receding horizon approach provides robustness to model errors and disturbances by continuously re-planning based on feedback.

\end{document}
